\section{Thảo luận \& Phân tích lỗi}

\subsection{Phân tích hiện tượng Overfitting/Underfitting}
% Ghi chú:
% - Dựa vào khoảng cách giữa kết quả Train và Test để kết luận trạng thái của mô hình.
% - Trình bày các phương pháp đã dùng để khắc phục: Early Stopping (dừng sớm khi eval_loss không giảm) và Weight Decay.

\subsection{Phân tích hiệu quả của mô hình lai (Hybrid Spam Model)}
% Ghi chú (MỚI THÊM):
% - Lấy ví dụ minh họa về 1-2 bình luận seeding tinh vi mà Rule-based (điểm = 0) bỏ qua nhưng Isolation Forest (score = -1) lại bắt được nhờ bất thường về TTR (Type-Token Ratio) hoặc độ dài từ.
% - Bàn luận về việc lọc Spam đã giúp làm sạch tập dữ liệu huấn luyện, từ đó nâng cao F1-Score của mô hình PhoBERT như thế nào.


\subsection{Phân tích các trường hợp sai với Explainable AI (LIME)}
% Ghi chú:
% - Lấy 1-2 ví dụ cụ thể về dữ liệu văn bản mà mô hình đoán sai.
% - Dùng kết quả từ thư viện LIME để đưa ra giả thuyết tại sao sai (Ví dụ: do khách hàng dùng icon mỉa mai, hoặc có từ khóa hiếm gặp gây nhiễu đặc trưng).


\subsection{So sánh hiệu năng}
% Ghi chú:
% - Lập bảng so sánh (Table) tổng hợp hiệu năng giữa các mô hình: Baseline (Auto-Weights), Baseline (SMOTE), và PhoBERT.
% - Đảm bảo bảng có tiêu đề bên trên và nhấn mạnh mô hình đạt kết quả Macro-F1 cao nhất.

\subsection{Phân tích lỗi (Error Analysis) trên Dữ liệu Đa phương thức}
% Ghi chú (CẬP NHẬT GỘP CHUNG TEXT VÀ IMAGE):
% - Phân tích Text (LIME): Lấy ví dụ bình luận mô hình đoán sai và dùng LIME giải thích.
% - Phân tích Image: Đánh giá chéo kết quả dự đoán của ResNet50 và CLIP. Trích xuất 1-2 hình ảnh thực tế mà ResNet50 đoán sai (Ví dụ: Góc chụp tối làm AI tưởng là hộp bị móp) để phân tích nguyên nhân.
% - Kết luận về sức mạnh của Cross-Modal Fusion: Chỉ ra rằng dù Text có thể đoán sai, hoặc Image đoán sai, nhưng khi kết hợp lại (Late Fusion) thành Trust Score, hệ thống vẫn đưa ra cảnh báo rủi ro chính xác cho người dùng.

\clearpage
